\section{Discussion}
\label{sec:discussion}

Our results move the hard part of writing parallel code \emph{twice}. We expected it to sit
in \emph{parallelization}---the races and order bugs a single run hides. For \Nmodels{}
models on standard OpenMP tasks, it did not: the models parallelized safely, and when they
were wrong the mistake was in the plain algorithm and already visible on one thread. The
bigger move is the one the speed measurement reveals: once correctness is mostly solved, the
hard part becomes \emph{parallelism itself}. The programs the check accepts run anywhere from
$0.33\times$ to $7.9\times$ (Section~\ref{sec:scaling}); a check that looks only at the output
cannot see the difference, because plain serial code maximizes it. Correctness stopped being
the hard part, and a check that scores only outputs is now measuring the wrong thing.

\vspace{3pt}
\noindent\textbf{The repair loop did the work, not the finer checker.}
Holding the model, prompt, and repair budget fixed and changing only the checking tool moves
robust correctness from \Lzerorate{} (no tool) to \Lonerate{} once there is a repair loop---
and then not at all: the full thread-sweep verifier reaches the same \Ltworate{}
($\Delta_{\text{L2--L1}}=\LtwoOverLone{}$), agreeing with the one-thread check on every one of
the \Npairs{} paired cells. The gain over single-shot is the loop's, not the checker's. This
does not contradict the design argument that a one-run check is \emph{structurally} blind to
nondeterministic races---that stays true---but it shows the argument only bites when the
models actually write such races, which here they did not.

\vspace{3pt}
\noindent\textbf{When the finer checker does pay off.}
The verifier is cheap---compile once, run a few thread counts a few times against the serial
reference---so its cost is a small, fixed premium. What our study measures is the
\emph{payout}, which is zero when the model$\times$task mix produces no silent races. It is
not zero in general: our injected-race control (Table~\ref{tab:injected}) confirms the
verifier rejects every racy program the one-thread check waves through. The payout grows with
how race-prone the mix is: weaker or non-code models that reach for manual shared-state updates
instead of a reduction; tasks whose natural parallelization invites a race; and the MPI and
CUDA settings our harness supports but we did not stress here. The right way to spend the
premium is to \emph{measure} it---run the differential gate, watch how often L1 and L2
disagree for your models and tasks, and keep it where they do.

\vspace{3pt}
\noindent\textbf{From a finding to a fix: a two-axis reward.}
The check an agent accepts on is also the reward a verifiable-reward RL loop would
optimize~\cite{rlvr,grpo}, and that is where the gameability stops being theoretical. Treat the
set of accepted programs for a task as the candidates a generator---or a best-of-$N$/GRPO
selection---chooses among. A correctness-only reward cannot choose: every robust program ties,
so an uninformed pick lands on the pool average. A two-axis reward (correct times fast) picks
the fastest. On the \NmultiTask{} tasks with more than one accepted program the gap is real: a
correctness-only reward realizes a mean \RewardRandom{}$\times$ 8-thread speedup against the
two-axis reward's \RewardTwoAxis{}$\times$ (\RewardGainPct{}\% more), and on the most spread-out
task it is the difference between a \RewardSpreadLo{}$\times$ and a \RewardSpreadHi{}$\times$
program the correctness reward calls equal. The fix, which we frame as future work, is a
\emph{second checked stage}: a performance gate that ends on a speed target and is
\emph{guarded} by the correctness gate---every optimization has to stay robust---so that chasing
the reward can never buy a wrong answer. This is the natural home for the negative result: once
correctness saturates, the reward should carry the second axis. We are careful to separate this
from prior work that rewards speed directly~\cite{pareval}; our point is narrower and about the
reward's \emph{specification}---a single-axis correctness reward is gameable by construction, and
we measure what that costs.

\vspace{3pt}
\noindent\textbf{What could be wrong.}
\emph{Construct.} We study shared-memory OpenMP over ParEval's tasks and serial references; the
MPI and CUDA cases, though our harness supports them, are left for future work, and the hazard
list is not exhaustive. Our main race evidence is output-differential, but we do not lean on it
alone: we cross-checked every accepted program with a happens-before detector (Archer/
ThreadSanitizer with OMPT annotations~\cite{archer,tsan}), which agreed---\NsanRaces{} races in
the \NsanTimed{} accepted programs---while correctly flagging our injected and open-model races.
So the null is backed by a second, independent instrument, not by one weak detector. We still
report it conservatively as a bound: with zero observed, the silent-failure risk is under about
$5\%$ at 95\% confidence (the rule of three). The speed finding does not share this fragility---
wall-time is measured directly.
\emph{Internal.} Each task has one serial reference; being sequential it is immune to the bugs
under test, but it does fix what the task means. Using one thread as the ``developer default''
is a modeling choice, picked as the most forgiving single run so the gap upper-bounds what a
lucky local check could hide.
\emph{External.} Our \Nmodels{} models come from one provider over \Ntasks{} tasks; holding
capability high sharpens the comparison but makes the near-absence of races unsurprising, and
neither the null nor the point estimates need carry to other vendors, weaker models, larger task
sets, or different prompt/temperature settings. That the null held even for a code-specialized
model softens, but does not remove, this. Because our harness reaches any model through one API,
cross-vendor and weaker-model replication is the most immediate extension. We report Wilson
intervals and read the point estimates alongside them, not as precise rankings; with
\Nmodels{}$\times$\Ntasks{} cells the intervals are wide.

\vspace{3pt}
\noindent\textbf{What to do about it.}
For teams adopting agentic assistants for HPC, a repair loop is clearly worth having---it gave
the only improvement we measured. A differential gate---check across thread counts and repeats
against a trusted reference, rather than trusting one good run---is cheap insurance we still
recommend, but with eyes open: for strong models on ordinary tasks its extra value over a
one-thread check may be nil, and the leftover failures are plain algorithm bugs that call for
better generation or repair, not finer checking. Treat the gate as instrumentation as much as
protection: it is also how you \emph{learn} whether your models write racy code.
